    \section*{Ejercicios del capítulo 13}

    \begin{ejercicio}
        Chebyshev probó que, si $\psi(x)/x$ tiende a un límite cuando $x\to\infty$ entonces este
        límite es igual a 1. Una demostración fue indicada ya en el ejercicio 4.26. Este
        ejercicio proporciona otra demostración basada en la identidad
        \begin{equation}\label{eq:13-37}
            -\frac{\zeta'(s)}{\zeta(s)}=s\int_1^{\infty}\frac{\psi(x)}{x^{s+1}}\,dx\qquad(\sigma>1)
        \end{equation}
        dada en el ejercicio 11.1(d).
        \begin{enumerate}
            \item[(a)] Probar que $(1-s)\zeta'(s)/\zeta(s)\to1$ cuando $s\to1+$.
            \item[(b)] Sea $\delta=\limsup_{x\to\infty}(\psi(x)/x)$. Dado $\varepsilon>0$,
            elegimos $N=N(\varepsilon)$ tal que $x\geq N$ implique
            $\psi(x)\leq(\delta+\varepsilon)x$. Tomamos $s$ real, $1<s\leq2$, descomponemos la
            integral de \eqref{eq:13-37} en dos partes, $\int_1^N+\int_N^{\infty}$, y acotamos
            cada parte a fin de obtener la desigualdad
            \begin{equation*}
                -\frac{\zeta'(s)}{\zeta(s)}\leq C(\varepsilon)+\frac{s(\delta+\varepsilon)}{s-1},
            \end{equation*}
            en donde $C(\varepsilon)$ es una constante independiente de $s$. Utilizar (a) para
            deducir que $\delta\geq1$.
            \item[(c)] Sea $\gamma=\liminf_{x\to\infty}(\psi(x)/x)$ y utilizar un razonamiento
            análogo para deducir que $\gamma\leq1$.
        \end{enumerate}
        Por consiguiente, si $\psi(x)/x$ tiende a un límite cuando $x\to\infty$, entonces
        $\gamma=\delta=1$.
    \end{ejercicio}

    \begin{ejercicio}
        Sea $A(x)=\sum_{n\leq x}a(n)$, en donde
        \begin{equation*}
            a(n)=
            \begin{cases}
                0&\text{si }n\neq\text{una potencia de primo,}\\
                \dfrac{1}{k}&\text{si }n=p^k.
            \end{cases}
        \end{equation*}
        Probar que $A(x)=\pi(x)+O(\sqrt{x}\log\log x)$.
    \end{ejercicio}

    \begin{ejercicio}
        \begin{enumerate}
            \item[(a)] Si $c>1$ y $x\neq$ entero, probar que, si $x>1$,
            \begin{equation*}
                \frac{1}{2\pi i}\int_{c-\infty i}^{c+\infty i}\log\zeta(s)\,\frac{x^s}{s}\,ds=\pi(x)+\frac{1}{2}\pi(x^{1/2})+\frac{1}{3}\pi(x^{1/3})+\cdots.
            \end{equation*}
            \item[(b)] Probar que el teorema del número primo es equivalente a la relación
            asintótica
            \begin{equation*}
                \frac{1}{2\pi i}\int_{c-\infty i}^{c+\infty i}\log\zeta(s)\,\frac{x^s}{s}\,ds\sim\frac{x}{\log x}\qquad\text{cuando }x\to\infty.
            \end{equation*}
        \end{enumerate}
        Una demostración del teorema del número primo basada en esta relación fue establecida
        por Landau en 1903.
    \end{ejercicio}

    \begin{ejercicio}
        Sea $M(x)=\sum_{n\leq x}\mu(n)$. Se desconoce el orden exacto de magnitud de $M(x)$ para
        $x$ grande. En el capítulo 4 probábamos que el teorema del número primo es equivalente a
        la relación $M(x)=o(x)$ cuando $x\to\infty$. Este ejercicio relaciona el orden de
        magnitud de $M(x)$ con la hipótesis de Riemann.

        Suponemos que existe una constante positiva $\theta$ tal que
        \begin{equation*}
            M(x)=O(x^{\theta})\qquad\text{para }x\geq1.
        \end{equation*}
        Probar que la fórmula
        \begin{equation*}
            \frac{1}{\zeta(s)}=s\int_1^{\infty}\frac{M(x)}{x^{s+1}}\,dx,
        \end{equation*}
        que es válida para $\sigma>1$ (ver ejercicio 11.1(c)) también es válida para
        $\sigma>\theta$. Deducir que $\zeta(s)\neq0$ para $\sigma>\theta$. En particular, esto
        prueba que la relación $M(x)=O(x^{1/2-\varepsilon})$ para cada $\varepsilon>0$ implica
        la hipótesis de Riemann. También se puede demostrar que la hipótesis de Riemann implica
        $M(x)=O(x^{1/2+\varepsilon})$ para cada $\varepsilon>0$. (Ver Titchmarsh
        \cite{titchmarsh1951}, p. 315.)
    \end{ejercicio}

    \begin{ejercicio}
        Probar el siguiente lema, que es análogo al lema \ref{lem:13-2}. Sea
        \begin{equation*}
            A_1(x)=\int_1^x\frac{A(u)}{u}\,du,
        \end{equation*}
        en donde $A(u)$ es una función creciente no negativa para $u\geq1$. Si, para un $c>0$ y
        un $L>0$, tenemos la fórmula asintótica
        \begin{equation*}
            A_1(x)\sim Lx^c\qquad\text{cuando }x\to\infty,
        \end{equation*}
        entonces es
        \begin{equation*}
            A(x)\sim cLx^c\qquad\text{cuando }x\to\infty.
        \end{equation*}
    \end{ejercicio}

    \begin{ejercicio}
        Probar que
        \begin{equation*}
            \frac{1}{2\pi i}\int_{2-\infty i}^{2+\infty i}\frac{y^s}{s^2}\,ds=0\qquad\text{si }0<y<1.
        \end{equation*}
        ¿Cuál es el valor de esta integral si $y\geq1$?
    \end{ejercicio}

    \begin{ejercicio}
        Expresar
        \begin{equation*}
            \frac{1}{2\pi i}\int_{2-\infty i}^{2+\infty i}\frac{x^s}{s^2}\left(-\frac{\zeta'(s)}{\zeta(s)}\right)ds
        \end{equation*}
        como una suma finita que contenga a $\Lambda(n)$.
    \end{ejercicio}

    \begin{ejercicio}
        Sea $\chi$ un carácter cualquiera de Dirichlet mód $k$ y $\chi_1$ el carácter principal.
        Definimos
        \begin{equation*}
            F(\sigma,t)=3\,\frac{L'}{L}(\sigma,\chi_1)+4\,\frac{L'}{L}(\sigma+it,\chi)+\frac{L'}{L}(\sigma+2it,\chi^2).
        \end{equation*}
        Si $\sigma>1$, probar que $F(\sigma,t)$ tiene parte real igual a
        \begin{equation*}
            -\sum_{n=1}^{\infty}\frac{\Lambda(n)}{n^{\sigma}}\operatorname{Re}\left\{3\chi_1(n)+4\chi(n)n^{-it}+\chi^2(n)n^{-2it}\right\}
        \end{equation*}
        y deducir que $\operatorname{Re}F(\sigma,t)\leq0$.
    \end{ejercicio}

    \begin{ejercicio}
        Suponemos que $L(s,\chi)$ tiene un cero de orden $m\geq1$ en $s=1+it$. Probar que, para
        este $t$, tenemos:
        \begin{enumerate}
            \item[(a)] $\displaystyle\frac{L'}{L}(\sigma+it,\chi)=\frac{m}{\sigma-1}+O(1)$ cuando $\sigma\to1+$,
        \end{enumerate}
        y
        \begin{enumerate}
            \item[(b)] existe un entero $r\geq0$ tal que
            \begin{equation*}
                \frac{L'}{L}(\sigma+2it,\chi^2)=\frac{r}{\sigma-1}+O(1)\qquad\text{cuando }\sigma\to1+,
            \end{equation*}
            excepto cuando $\chi^2=\chi_1$ y $t=0$.
        \end{enumerate}
    \end{ejercicio}

    \begin{ejercicio}
        Utilizar los ejercicios 8 y 9 para demostrar que
        \begin{equation*}
            L(1+it,\chi)\neq0\qquad\text{para todo real }t\text{ si }\chi^2\neq\chi_1
        \end{equation*}
        y que
        \begin{equation*}
            L(1+it,\chi)\neq0\qquad\text{para todo real }t\neq0\text{ si }\chi^2=\chi_1.
        \end{equation*}
        [\textit{Indicación:} Considerar $F(\sigma,t)$ cuando $\sigma\to1+$.]
    \end{ejercicio}

    \begin{ejercicio}
        Para toda función aritmética $f(n)$, probar que las siguientes afirmaciones son
        equivalentes:
        \begin{enumerate}
            \item[(a)] $f(n)=O(n^{\varepsilon})$ para cada $\varepsilon>0$ y todo
            $n\geq n_1$.
            \item[(b)] $f(n)=o(n^{\delta})$ para cada $\delta>0$ cuando $n\to\infty$.
        \end{enumerate}
    \end{ejercicio}

    \begin{ejercicio}
        Sea $f$ una función multiplicativa tal que, si $p$ es primo, entonces
        \begin{equation*}
            f(p^m)\to0\qquad\text{cuando }p^m\to\infty.
        \end{equation*}
        Es decir, para cada $\varepsilon>0$, existe un $N(\varepsilon)$ tal que
        $\abs{f(p^m)}<\varepsilon$ si $p^m>N(\varepsilon)$. Probar que $f(n)\to0$ cuando
        $n\to\infty$.

        [\textit{Indicación:} Existe una constante $A>0$ tal que $\abs{f(p^m)}<A$ para todo
        primo $p$ y todo $m>0$, y una constante $B>0$ tal que $\abs{f(p^m)}<1$ si $p^m>B$.]
    \end{ejercicio}

    \begin{ejercicio}
        Si $\alpha\geq0$, sea $\sigma_{\alpha}(n)=\sum_{d\mid n}d^{\alpha}$. Probar que, para
        cada $\delta>0$, tenemos
        \begin{equation*}
            \sigma_{\alpha}(n)=o(n^{\alpha+\delta})\qquad\text{cuando }n\to\infty.
        \end{equation*}
        [\textit{Indicación:} Utilizar el ejercicio 12.]
    \end{ejercicio}

